\documentclass[11pt]{article}
\usepackage[utf8]{inputenc}
\usepackage[T1]{fontenc}
\usepackage{amsmath,amssymb,amsthm}
\usepackage{graphicx}
\usepackage{booktabs}
\usepackage{listings}
\usepackage{xcolor}
\usepackage{subcaption}
\usepackage[margin=2.5cm]{geometry}
\usepackage{hyperref}

\graphicspath{{../results/experiments/}{../results/benchmark/}}

\newtheorem{proposition}{Proposition}
\DeclareMathOperator{\Tr}{Tr}
\newcommand{\C}{\mathbb{C}}
\newcommand{\M}{\mathcal{M}}

\lstset{language=Python, basicstyle=\ttfamily\small, numbers=left, numberstyle=\tiny,
        frame=single, breaklines=true, columns=fullflexible, keepspaces=true}

\title{Finite-Dimensional Type I von Neumann Algebras in PyTorch:\\
A GPU-Accelerated Framework for Random Block-Diagonal Operators}
\author{Irina Nikolaeva \and Andrej Novikov}
\date{}

\begin{document}
\maketitle

\begin{abstract}
We present \texttt{torch\_vn\_algebra}, an open-source Python library built on PyTorch for
numerical experiments with finite-dimensional Type~I von Neumann algebras (direct sums of matrix
algebras). The library provides:
\begin{itemize}
\item a compact batched tensor representation $(B, C, k_{\max}, k_{\max})$ that handles both Monte
      Carlo samples and multiple direct summands;
\item lazy evaluation of operators to avoid unnecessary memory allocation;
\item generation of random operators with arbitrary eigenvalue distributions (user-provided
      samplers) and several unitary ensembles (Haar on $U(n)$ and $O(n)$, $SU(n)$ and $SO(n)$,
      COE, CSE, diagonal phases);
\item functional calculus via SVD (absolute value, square root, inverse, entropy) and a hybrid
      method for extreme eigenvalues (exact diagonalisation for blocks of size $k_c\le 256$,
      otherwise shifted power iteration);
\item three trace functionals (blunt, normalised subspace trace, and the von Neumann tracial
      state);
\item GPU-accelerated batched linear algebra for moderate-scale Monte Carlo studies (e.g.\
      $2\times 10^4$ samples of operators with up to 32 blocks).
\end{itemize}
The library is validated against analytical expectations (Haar moments, trace properties).
Performance benchmarks on a Tesla P100 GPU are presented and discussed, and three Monte Carlo
experiments on trace inequalities are reported. Limitations and future work are outlined. The
code is open-source.
\end{abstract}

\section{Introduction}

Finite-dimensional Type I von Neumann algebras,
\[
  \M = \bigoplus_{c=1}^{C} M_{n_c}(\C),
\]
naturally appear in quantum mechanics with superselection rules, decoherence, and in random
matrix theory when symmetries or constraints force operators into block-diagonal form. Monte
Carlo simulations of such systems -- averaging over many random operators with controlled
spectral properties -- require efficient numerical tools that can exploit modern hardware.
General-purpose libraries (NumPy/SciPy, QuTiP, Qiskit) do not provide a native batched
representation of block-diagonal (direct-sum) operators together with arbitrary prescribed
spectra and GPU execution.

\texttt{torch\_vn\_algebra} fills this gap by combining a compact tensor representation that
batches both Monte Carlo samples and direct summands with padding to a common $k_{\max}$; lazy
evaluation (operators are built from generators and materialised only when needed); flexible
random operators (any eigenvalue distribution via a user callable) with several unitary
ensembles; batched functional calculus with hybrid eigenvalue extraction; three trace
functionals $\Tr_{\mathrm{blunt}}$, $\Tr_{\mathrm{norm}}$ and the von Neumann tracial state
$\tau_{\mathrm{vN}}$; and full integration with PyTorch's GPU backend.

The library builds upon the theoretical framework developed in~\cite{abed2024,novikov2020},
where the Michelson contrast was generalised to operators in von Neumann algebras. The present
work provides a scalable GPU implementation of these concepts for Monte Carlo studies.

The paper is organised as follows. Section~\ref{sec:background} recalls necessary definitions.
Section~\ref{sec:implementation} describes the implementation. Section~\ref{sec:benchmarks}
presents validation and performance benchmarks. Section~\ref{sec:experiments} gives Monte Carlo
experiments. Section~\ref{sec:limitations} discusses limitations and future work.
Section~\ref{sec:conclusion} concludes.

\section{Mathematical background}\label{sec:background}

A finite-dimensional Type I von Neumann algebra is $\M = \bigoplus_{c=1}^{C} M_{n_c}(\C)$ acting
on $\mathcal H = \bigoplus_{c=1}^{C} \C^{k_c}$ with $k_c \le n_c$. Operators are block-diagonal:
$A = \bigoplus_c A_c$. Three trace functionals are relevant:
\[
  \Tr_{\mathrm{blunt}}(A) = \sum_{c=1}^{C} \Tr(A_c), \qquad
  \Tr_{\mathrm{norm}}(A) = \sum_{c=1}^{C} \frac{1}{k_c}\Tr\bigl(A_c^{(k_c)}\bigr), \qquad
  \tau_{\mathrm{vN}}(A) = \frac{1}{C}\sum_{c=1}^{C} \frac{1}{k_c}\Tr\bigl(A_c^{(k_c)}\bigr),
\]
where $A_c^{(k_c)}$ is the restriction of $A_c$ to the $k_c$-dimensional subspace.
$\tau_{\mathrm{vN}}$ is a faithful normal tracial state. For $A\in\M$ we write
$|A| = (A^*A)^{1/2}$. For a density matrix $\rho$ ($\rho\ge 0$, $\Tr\rho = 1$), the von Neumann
entropy is $S(\rho) = -\Tr(\rho\log\rho)$. For a positive operator $A$, define its Michelson
contrast as
\[
  \Delta(A) = \frac{\lambda_{\max}(A) - \lambda_{\min}(A)}{\lambda_{\max}(A) + \lambda_{\min}(A)},
\]
where the extreme eigenvalues are taken over the whole algebra, i.e.\ over all blocks.
$\Delta(A)\in[0,1]$, and $\Delta(A) = 0$ if and only if $A$ is a scalar multiple of the
identity.

The theory of noncommutative integration associates to a von Neumann algebra $\M$ equipped with
a faithful normal trace $\tau$ a family of noncommutative $L^p$-spaces. In the
finite-dimensional Type I case, the $L^1$-space is simply $\M$ with the norm
$\|A\|_1 = \tau(|A|)$. For a positive operator $X$ affiliated with $\M$, one can define its
$L^1$-norm as $\tau(X)$ when finite~\cite{novikov2017}.

\section{Implementation}\label{sec:implementation}

Operators are stored as a 4-D tensor $(B, C, k_{\max}, k_{\max})$ with the active block of size
$k_c\times k_c$ in the top-left corner; the remaining entries are zero padding. The
\texttt{Operator} class uses lazy generation: the matrix is materialised only when accessed.
The auxiliary class \texttt{HilbertSpace} provides inner product, norm, and orthonormal bases
(standard, random, Haar), supporting embedding and restriction of vectors and operators.

\paragraph{Random operators.} Random operator generation is performed via
\texttt{operator\_from\_eigenvalues}. The user supplies a callable
\texttt{eigenvalue\_sampler(dim)} returning either one spectrum of length \texttt{dim} or a
batch of spectra of shape $(B, \mathtt{dim})$. For each channel a batch of random unitaries $U$ is
drawn from the chosen ensemble and $A_c = U\,\mathrm{diag}(\lambda)\,U^*$ is formed (and
symmetrised when the spectrum is real). Properties (self-adjointness, positivity, etc.) are
inferred from the spectrum; forced properties are validated when the matrix is materialised,
with tolerances relative to the working precision.

\paragraph{Unitary ensembles.} Haar-distributed unitaries are sampled by the QR decomposition of
a Ginibre matrix with i.i.d.\ circular complex (or real) standard Gaussian entries, followed by
fixing the phases of the diagonal of $R$~\cite{mezzadri2007}. $SU(n)$ is obtained by multiplying a
Haar unitary by $\det(U)^{-1/n}$ and $SO(n)$ by flipping the sign of one column when
$\det = -1$; both maps push Haar measure forward to Haar measure. The circular orthogonal and
symplectic ensembles are $S = W^{T} W$ and $S = W^{R} W$, $W^{R} = J W^{T} J^{T}$, with $W$ Haar
on $U(n)$~\cite{mehta2004}. All samplers are batched.

\paragraph{Extreme eigenvalues.} For self-adjoint operators, $\lambda_{\max}$ and
$\lambda_{\min}$ are obtained by exact diagonalisation (\texttt{torch.linalg.eigvalsh}) when all
blocks satisfy $k_c\le 256$ (threshold \texttt{exact\_eig\_max\_dim}). For larger blocks a
two-stage power iteration is used in double precision: first the eigenvalue $\mu$ of largest
modulus is found, then the dominant eigenvalue of $\mu I - A$, which gives the opposite end of
the spectrum (at most 5000 iterations, stopping when the Rayleigh quotient changes by less than
$10^{-12}$). The operator norm uses a shorter single-precision power iteration on $A^*A$ (at
most 100 iterations, tolerance $10^{-8}$).

\paragraph{Functional calculus and traces.} $|A|$, $\sqrt{A}$ (for $A\ge 0$), the
(pseudo-)inverse $A^{-1}$ and $\Tr(A\log A)$ are computed via batched SVD,
$A = W\Sigma V^*$, $|A| = V\Sigma V^*$. The three traces are implemented as
\texttt{Tr\_blunt}, \texttt{Tr\_norm} and \texttt{tau\_vN}; the Michelson contrast
\texttt{michelson\_contrast} returns $\Delta(A)$ for every sample of the batch.

\section{Validation and performance benchmarks}\label{sec:benchmarks}

\subsection{Validation}

All validation runs are reproduced by \texttt{scripts/validation.py}.

\paragraph{Haar measure.} For $U$ Haar-distributed on $U(n)$ one has
$\mathbb E|U_{11}|^2 = 1/n$ and, by invariance under $U\mapsto e^{i\varphi}U$,
$\mathbb E[U_{11}^2] = 0$. The second identity fails for non-circular Gaussian input and is
therefore a useful check of the sampler. Table~\ref{tab:haar} shows the results for
$n = 2,4,8,16,32$ using $10\,000$ samples generated by the library. All relative errors are
within about two standard errors of the mean.

\begin{table}[h]
\centering
\begin{tabular}{rrrrr}
\toprule
$n$ & $1/n$ & relative error of $\mathbb E|U_{11}|^2$ & relative standard error & $|\mathbb E[U_{11}^2]|$\\
\midrule
2  & 0.5     & $9.3\times10^{-4}$ & $5.8\times10^{-3}$ & $3.3\times10^{-3}$\\
4  & 0.25    & $5.1\times10^{-3}$ & $7.8\times10^{-3}$ & $3.5\times10^{-3}$\\
8  & 0.125   & $4.7\times10^{-3}$ & $8.9\times10^{-3}$ & $6.6\times10^{-4}$\\
16 & 0.0625  & $1.9\times10^{-2}$ & $9.4\times10^{-3}$ & $3.0\times10^{-4}$\\
32 & 0.03125 & $2.0\times10^{-2}$ & $9.6\times10^{-3}$ & $2.7\times10^{-4}$\\
\bottomrule
\end{tabular}
\caption{Haar measure on $U(n)$, $10\,000$ samples.}
\label{tab:haar}
\end{table}

\paragraph{Power iteration.} For a $3\times3$ matrix with eigenvalues $[1.0, 0.99, 0.98]$
(gap $0.01$), the power iteration needed $699$ iterations before successive Rayleigh quotients
differed by less than $10^{-8}$, compared to $15$ iterations for the well-separated spectrum
$[2, 1, 0.5]$. Note that for the small gap the actual error at that point was
$4.9\times10^{-7}$, i.e.\ a stopping rule on successive differences overestimates the accuracy
when the gap is small. This is why the library uses exact diagonalisation up to $k_c = 256$.

\paragraph{Square root.} We generated $1000$ random $100\times 100$ positive matrices
(eigenvalues uniform in $[0,1]$) and computed $B = \sqrt{A}$ via SVD. The relative error
$\|B^2 - A\|_F/\|A\|_F$ was $2.0\times10^{-6}$ on average (maximum $2.4\times10^{-6}$) in the
library's default single precision and $4.0\times10^{-15}$ (maximum $4.9\times10^{-15}$) in double
precision, i.e.\ at the level of the respective machine epsilon.

\subsection{Benchmark setup and speedup results}

All benchmarks were run on an NVIDIA Tesla P100 (16\,GB) and an Intel Xeon E5-2650 v4
@ 2.2\,GHz (12 cores) with \texttt{scripts/benchmark.py}. Each test input is a tensor of shape
$(C, k, k)$, i.e.\ the channel axis plays the role of the batch, with
$k\in\{2,4,\dots,64\}$ and $C\in\{1,2,4,\dots,1024\}$; every configuration was timed with three
random seeds and ten repetitions. Speedup is defined as $t_{\mathrm{CPU}}/t_{\mathrm{GPU}}$ and
includes the final device synchronisation but not the host-to-device transfer. The benchmark
times the batched PyTorch kernels underlying the library operations (trace, addition,
multiplication, inverse, SVD, symmetric eigensolver, norms) on Gaussian random matrices; the
CPU used PyTorch's default thread count.
% TODO(authors): confirm the CPU thread setting of the stored run.

Table~\ref{tab:speedup} shows the speedup for the inverse and for matrix multiplication. A
second, independent run with the same configuration (stored in the repository as
\texttt{results/benchmark/run\_b}) agrees with the first within a median factor of $1.12$, but
single configurations may differ considerably, so the numbers should be read as orders of
magnitude. For a single channel the CPU is always faster (speedup $0.3$--$0.7$), because kernel
launch overhead dominates. The GPU pays off from about $C\ge 64$ blocks of size $k\ge 8$, with
speedups of 15--135 at $C = 1024$. For the SVD and the symmetric eigensolver the GPU becomes
\emph{slower} than the CPU at $k = 64$ (speedup $0.2$--$0.6$ for all $C$), which marks the
transition from the batched small-matrix GPU solvers to the general ones; this is the main
practical limitation for larger blocks (Section~\ref{sec:limitations}).

\begin{table}[h]
\centering
\begin{tabular}{r|rrr|rrr}
\toprule
 & \multicolumn{3}{c|}{inverse} & \multicolumn{3}{c}{multiplication}\\
$k$ & $C=1$ & $C=64$ & $C=1024$ & $C=1$ & $C=64$ & $C=1024$\\
\midrule
2  & 0.3 & 1.6 (1.0)   & 11.3  & 0.4 & 0.5 (0.5)  & 0.5\\
4  & 0.3 & 1.5 (1.2)   & 15.5  & 0.4 & 0.5 (0.4)  & 1.0\\
8  & 0.3 & 25.1 (16.0) & 134.8 & 0.5 & 2.5 (2.0)  & 27.2\\
16 & 0.4 & 21.0 (19.1) & 50.6  & 0.5 & 2.8 (2.2)  & 34.8\\
32 & 0.3 & 34.7 (40.5) & 49.0  & 0.6 & 6.0 (4.1)  & 58.3\\
64 & 0.5 & 17.9 (10.7) & 18.9  & 0.7 & 11.0 (10.6) & 138.4\\
\bottomrule
\end{tabular}
\caption{Mean speedup $t_{\mathrm{CPU}}/t_{\mathrm{GPU}}$ for $C$ blocks of size $k$
(first run; values of the second run in parentheses).}
\label{tab:speedup}
\end{table}

\section{Monte Carlo experiments on trace inequalities}\label{sec:experiments}

We perform three Monte Carlo experiments that mirror earlier NumPy code and illustrate the
library's capabilities. All experiments are reproduced by \texttt{scripts/experiment.py}. For
each block size $k\in\{2,16\}$ and number of channels $C\in\{1,2,16,32\}$ (all $k_c = k$) we
generate $20\,000$ random operator pairs in the real algebra
$\M = \bigoplus_{c=1}^C M_k(\mathbb R)$. Operators are constructed via the spectral theorem: in
every block the largest eigenvalue is $1$, the smallest is $(1-\delta)/(1+\delta)$ with $\delta$
uniform in $[0,1)$, and the remaining eigenvalues are uniform in between; the eigenvectors are
Haar-random orthogonal matrices, independent for every block and sample. Throughout this
section $\Tr = \Tr_{\mathrm{blunt}}$, and ``$U$ unitary in $\M$'' means
$U = \bigoplus_c U_c$ with independent Haar-random orthogonal $U_c$.

Table~\ref{tab:experiments} summarises the sign of the quantity $z$ in each experiment; the
figures show scatter plots of $z$ against the contrasts. Raw samples, all plots, and the
summary table are available in the repository (\texttt{results/experiments}).

\begin{table}[h]
\centering
\begin{tabular}{rr|rr|rr|rr}
\toprule
 & & \multicolumn{2}{c|}{Exp.~1} & \multicolumn{2}{c|}{Exp.~2} & \multicolumn{2}{c}{Exp.~3}\\
$k$ & $C$ & $z>0$, \% & $\max z$ & $z>0$, \% & $z<0$, \% & $z<0$, \% & $\min z$\\
\midrule
2 & 1 & 7.0 & 0.227 & 52.9 & 46.6 & 0.0 & -8.34e-07 \\
2 & 2 & 0.7 & 0.328 & 54.2 & 45.8 & 0.0 & -9.54e-07 \\
2 & 16 & 0.0 & -4.37 & 62.4 & 37.6 & 0.0 & 0.00115 \\
2 & 32 & 0.0 & -15.9 & 68.2 & 31.8 & 0.0 & 0.0764 \\
16 & 1 & 0.0 & -2.18 & 57.4 & 41.9 & 0.0 & -7.63e-06 \\
16 & 2 & 0.0 & -5.73 & 63.0 & 36.9 & 0.0 & 1.72e-05 \\
16 & 16 & 0.0 & -89 & 86.4 & 13.6 & 0.0 & 1.83 \\
16 & 32 & 0.0 & -193 & 94.1 & 5.9 & 0.0 & 6.11 \\
\bottomrule
\end{tabular}
\caption{Sign of $z$ over $20\,000$ samples per configuration ($|z|\le10^{-4}$ counted as zero).}
\label{tab:experiments}
\end{table}

\subsection{Experiment 1 (positive $X$, positive $Y$, unitary $U$)}

Both $X$ and $Y$ are positive operators and $U$ is a random unitary in $\M$. The quantity of
interest is
\[
  z = |\Tr(XUY)| - \Tr(XY).
\]
The inequality $|\Tr(XUY)|\le\Tr(XY)$ holds for all unitaries $U$ when $X$ and $Y$ commute, but
not in general: writing the polar decomposition $YX = W|YX|$ with $W$ unitary in $\M$ and
choosing $U = W^*$ gives $\Tr(XUY) = \Tr(W^*YX) = \|YX\|_1 \ge \Tr(YX) = \Tr(XY)$, with strict
inequality whenever $YX$ is not positive. Random unitaries come close to this extremal choice only rarely, and
increasingly so as $k$ and $C$ grow. In the simulations $z>0$ occurs in $7.0\%$ of the samples
for $k = 2$, $C = 1$ and in $0.7\%$ for $k=2$, $C=2$, but in none of the samples with $k = 16$
or $C\ge 16$. Positive values occur only away from scalar operators: for $k=2$, $C=1$ the
largest $z$ over samples with $\Delta(X),\Delta(Y)<0.1$ is $-4\times10^{-4}$, and the size of the
violations grows with $\Delta(X)$ (Figure~\ref{fig:exp1-small}). The contrasts $\Delta(X)$ and $\Delta(Y)$ are computed from the global extreme
eigenvalues.

\begin{lstlisting}[caption={Experiment 1 implementation}]
def experiment1(alg, batch_size, num_batches, device):
    pos = positive_sampler(batch_size, device)
    out = []
    for _ in range(num_batches):
        X = alg.operator_from_eigenvalues(pos, batch_size=batch_size,
                                          force_positive=True, force_self_adjoint=True)
        Y = alg.operator_from_eigenvalues(pos, batch_size=batch_size,
                                          force_positive=True, force_self_adjoint=True)
        U = alg.random_unitary_operator(batch_size)
        z = (X @ U @ Y).trace.abs() - (X @ Y).trace.real
        out.append((X.michelson_contrast, Y.michelson_contrast, z))
    return collect(out)
\end{lstlisting}

\begin{figure}[p]
\centering
\begin{subfigure}{\textwidth}\centering
\includegraphics[width=0.32\textwidth]{exp1/plots/exp1_dim2_ch1_xz.png}
\includegraphics[width=0.32\textwidth]{exp1/plots/exp1_dim2_ch1_yz.png}
\includegraphics[width=0.32\textwidth]{exp1/plots/exp1_dim2_ch1_3d.png}
\caption{$k=2$, $C=1$}
\end{subfigure}
\begin{subfigure}{\textwidth}\centering
\includegraphics[width=0.32\textwidth]{exp1/plots/exp1_dim2_ch2_xz.png}
\includegraphics[width=0.32\textwidth]{exp1/plots/exp1_dim2_ch2_yz.png}
\includegraphics[width=0.32\textwidth]{exp1/plots/exp1_dim2_ch2_3d.png}
\caption{$k=2$, $C=2$}
\end{subfigure}
\begin{subfigure}{\textwidth}\centering
\includegraphics[width=0.32\textwidth]{exp1/plots/exp1_dim2_ch16_xz.png}
\includegraphics[width=0.32\textwidth]{exp1/plots/exp1_dim2_ch16_yz.png}
\includegraphics[width=0.32\textwidth]{exp1/plots/exp1_dim2_ch16_3d.png}
\caption{$k=2$, $C=16$}
\end{subfigure}
\caption{Experiment 1: small dimension ($k = 2$). Left: $z$ vs.\ $\Delta(X)$; middle: $z$ vs.\ $\Delta(Y)$; right: both.}
\label{fig:exp1-small}
\end{figure}

\begin{figure}[p]
\centering
\begin{subfigure}{\textwidth}\centering
\includegraphics[width=0.32\textwidth]{exp1/plots/exp1_dim16_ch1_xz.png}
\includegraphics[width=0.32\textwidth]{exp1/plots/exp1_dim16_ch1_yz.png}
\includegraphics[width=0.32\textwidth]{exp1/plots/exp1_dim16_ch1_3d.png}
\caption{$k=16$, $C=1$}
\end{subfigure}
\begin{subfigure}{\textwidth}\centering
\includegraphics[width=0.32\textwidth]{exp1/plots/exp1_dim16_ch16_xz.png}
\includegraphics[width=0.32\textwidth]{exp1/plots/exp1_dim16_ch16_yz.png}
\includegraphics[width=0.32\textwidth]{exp1/plots/exp1_dim16_ch16_3d.png}
\caption{$k=16$, $C=16$}
\end{subfigure}
\begin{subfigure}{\textwidth}\centering
\includegraphics[width=0.32\textwidth]{exp1/plots/exp1_dim16_ch32_xz.png}
\includegraphics[width=0.32\textwidth]{exp1/plots/exp1_dim16_ch32_yz.png}
\includegraphics[width=0.32\textwidth]{exp1/plots/exp1_dim16_ch32_3d.png}
\caption{$k=16$, $C=32$}
\end{subfigure}
\caption{Experiment 1: intermediate dimension ($k = 16$). Left: $z$ vs.\ $\Delta(X)$; middle: $z$ vs.\ $\Delta(Y)$; right: both.}
\end{figure}


\subsection{Experiment 2 (non-Hermitian $X$, non-Hermitian $Y$)}

Both $X$ and $Y$ are made non-Hermitian by multiplying positive operators with random unitaries:
$X = UX_0$, $Y = VY_0$, so that $|X| = X_0$ and $|Y| = Y_0$. The quantity of interest is
\[
  z = \Tr|XY| - \Tr(|X|\,|Y|) = \|X_0 V Y_0\|_1 - \Tr(X_0 Y_0).
\]
$z$ takes both signs: for $V = I$ one has $z = \|X_0Y_0\|_1 - \Tr(X_0Y_0)\ge 0$, while for $V$
far from the identity $\|X_0VY_0\|_1$ can be much smaller than $\Tr(X_0Y_0)$. In the
simulations $z>0$ in $53$--$94\%$ of the samples, the share growing with $k$ and $C$. The contrasts are
computed from $|X| = X_0$ and $|Y| = Y_0$, i.e.\ from the singular values of $X$ and $Y$.

\begin{lstlisting}[caption={Experiment 2 implementation}]
def experiment2(alg, batch_size, num_batches, device):
    pos = positive_sampler(batch_size, device)
    out = []
    for _ in range(num_batches):
        X0 = alg.operator_from_eigenvalues(pos, batch_size=batch_size,
                                           force_positive=True, force_self_adjoint=True)
        Y0 = alg.operator_from_eigenvalues(pos, batch_size=batch_size,
                                           force_positive=True, force_self_adjoint=True)
        X = alg.random_unitary_operator(batch_size) @ X0
        Y = alg.random_unitary_operator(batch_size) @ Y0
        z = (X @ Y).abs().trace.real - (X.abs() @ Y.abs()).trace.real
        out.append((X0.michelson_contrast, Y0.michelson_contrast, z))
    return collect(out)
\end{lstlisting}

\begin{figure}[p]
\centering
\begin{subfigure}{\textwidth}\centering
\includegraphics[width=0.32\textwidth]{exp2/plots/exp2_dim2_ch1_xz.png}
\includegraphics[width=0.32\textwidth]{exp2/plots/exp2_dim2_ch1_yz.png}
\includegraphics[width=0.32\textwidth]{exp2/plots/exp2_dim2_ch1_3d.png}
\caption{$k=2$, $C=1$}
\end{subfigure}
\begin{subfigure}{\textwidth}\centering
\includegraphics[width=0.32\textwidth]{exp2/plots/exp2_dim2_ch2_xz.png}
\includegraphics[width=0.32\textwidth]{exp2/plots/exp2_dim2_ch2_yz.png}
\includegraphics[width=0.32\textwidth]{exp2/plots/exp2_dim2_ch2_3d.png}
\caption{$k=2$, $C=2$}
\end{subfigure}
\begin{subfigure}{\textwidth}\centering
\includegraphics[width=0.32\textwidth]{exp2/plots/exp2_dim2_ch16_xz.png}
\includegraphics[width=0.32\textwidth]{exp2/plots/exp2_dim2_ch16_yz.png}
\includegraphics[width=0.32\textwidth]{exp2/plots/exp2_dim2_ch16_3d.png}
\caption{$k=2$, $C=16$}
\end{subfigure}
\caption{Experiment 2: small dimension ($k = 2$). Left: $z$ vs.\ $\Delta(|X|)$; middle: $z$ vs.\ $\Delta(|Y|)$; right: both.}
\end{figure}

\begin{figure}[p]
\centering
\begin{subfigure}{\textwidth}\centering
\includegraphics[width=0.32\textwidth]{exp2/plots/exp2_dim16_ch1_xz.png}
\includegraphics[width=0.32\textwidth]{exp2/plots/exp2_dim16_ch1_yz.png}
\includegraphics[width=0.32\textwidth]{exp2/plots/exp2_dim16_ch1_3d.png}
\caption{$k=16$, $C=1$}
\end{subfigure}
\begin{subfigure}{\textwidth}\centering
\includegraphics[width=0.32\textwidth]{exp2/plots/exp2_dim16_ch16_xz.png}
\includegraphics[width=0.32\textwidth]{exp2/plots/exp2_dim16_ch16_yz.png}
\includegraphics[width=0.32\textwidth]{exp2/plots/exp2_dim16_ch16_3d.png}
\caption{$k=16$, $C=16$}
\end{subfigure}
\begin{subfigure}{\textwidth}\centering
\includegraphics[width=0.32\textwidth]{exp2/plots/exp2_dim16_ch32_xz.png}
\includegraphics[width=0.32\textwidth]{exp2/plots/exp2_dim16_ch32_yz.png}
\includegraphics[width=0.32\textwidth]{exp2/plots/exp2_dim16_ch32_3d.png}
\caption{$k=16$, $C=32$}
\end{subfigure}
\caption{Experiment 2: intermediate dimension ($k = 16$). Left: $z$ vs.\ $\Delta(|X|)$; middle: $z$ vs.\ $\Delta(|Y|)$; right: both.}
\end{figure}


\subsection{Experiment 3 (self-adjoint $X$, positive $Y$)}

Here $X$ is self-adjoint (eigenvalues with independent random signs) and $Y$ is positive. The
quantity of interest is
\[
  z = \Tr(Y|X|Y) - \Tr|YXY|.
\]

\begin{proposition}
For $X = X^*$ and $Y\ge 0$ in $\M$ one has $z\ge 0$. If $Y$ commutes with $X$ (in particular,
if $Y$ is central), then $z = 0$.
\end{proposition}
\begin{proof}
Write $X = X_+ - X_-$ with $X_\pm\ge 0$, $|X| = X_+ + X_-$. By the triangle inequality for the
trace norm, $\Tr|YXY| = \|YX_+Y - YX_-Y\|_1 \le \Tr(YX_+Y) + \Tr(YX_-Y) = \Tr(Y|X|Y)$. If
$YX = XY$, then $YXY = Y^2X$ and $|YXY| = Y^2|X| = Y|X|Y$.
\end{proof}

The simulations agree with the proposition: no sample has $z<0$ beyond round-off
($|z|\le 10^{-5}$). $z = 0$ also occurs whenever all eigenvalues of $X$ have the same sign, which
for $k = 2$, $C = 1$ happens in half of the samples. The size of $z$ thus measures how far $Y$
is from commuting with $X$; its dependence on $\Delta(Y)$ is shown in the figures. The relation of
this quantity to the characterisation of central elements by trace inequalities~\cite{novikov2015}
and to noncommutative $L^1$-spaces~\cite{novikov2017} is a subject for further study.

\begin{lstlisting}[caption={Experiment 3 implementation}]
def experiment3(alg, batch_size, num_batches, device):
    sa = selfadjoint_sampler(batch_size, device)
    pos = positive_sampler(batch_size, device)
    out = []
    for _ in range(num_batches):
        X = alg.operator_from_eigenvalues(sa, batch_size=batch_size, force_self_adjoint=True)
        Y = alg.operator_from_eigenvalues(pos, batch_size=batch_size,
                                          force_positive=True, force_self_adjoint=True)
        absX = X.abs()
        z = (Y @ absX @ Y).trace.real - (Y @ X @ Y).abs().trace.real
        out.append((absX.michelson_contrast, Y.michelson_contrast, z))
    return collect(out)
\end{lstlisting}

\begin{figure}[p]
\centering
\begin{subfigure}{\textwidth}\centering
\includegraphics[width=0.32\textwidth]{exp3/plots/exp3_dim2_ch1_xz.png}
\includegraphics[width=0.32\textwidth]{exp3/plots/exp3_dim2_ch1_yz.png}
\includegraphics[width=0.32\textwidth]{exp3/plots/exp3_dim2_ch1_3d.png}
\caption{$k=2$, $C=1$}
\end{subfigure}
\begin{subfigure}{\textwidth}\centering
\includegraphics[width=0.32\textwidth]{exp3/plots/exp3_dim2_ch2_xz.png}
\includegraphics[width=0.32\textwidth]{exp3/plots/exp3_dim2_ch2_yz.png}
\includegraphics[width=0.32\textwidth]{exp3/plots/exp3_dim2_ch2_3d.png}
\caption{$k=2$, $C=2$}
\end{subfigure}
\begin{subfigure}{\textwidth}\centering
\includegraphics[width=0.32\textwidth]{exp3/plots/exp3_dim2_ch16_xz.png}
\includegraphics[width=0.32\textwidth]{exp3/plots/exp3_dim2_ch16_yz.png}
\includegraphics[width=0.32\textwidth]{exp3/plots/exp3_dim2_ch16_3d.png}
\caption{$k=2$, $C=16$}
\end{subfigure}
\caption{Experiment 3: small dimension ($k = 2$). Left: $z$ vs.\ $\Delta(|X|)$; middle: $z$ vs.\ $\Delta(Y)$; right: both.}
\end{figure}

\begin{figure}[p]
\centering
\begin{subfigure}{\textwidth}\centering
\includegraphics[width=0.32\textwidth]{exp3/plots/exp3_dim16_ch1_xz.png}
\includegraphics[width=0.32\textwidth]{exp3/plots/exp3_dim16_ch1_yz.png}
\includegraphics[width=0.32\textwidth]{exp3/plots/exp3_dim16_ch1_3d.png}
\caption{$k=16$, $C=1$}
\end{subfigure}
\begin{subfigure}{\textwidth}\centering
\includegraphics[width=0.32\textwidth]{exp3/plots/exp3_dim16_ch16_xz.png}
\includegraphics[width=0.32\textwidth]{exp3/plots/exp3_dim16_ch16_yz.png}
\includegraphics[width=0.32\textwidth]{exp3/plots/exp3_dim16_ch16_3d.png}
\caption{$k=16$, $C=16$}
\end{subfigure}
\begin{subfigure}{\textwidth}\centering
\includegraphics[width=0.32\textwidth]{exp3/plots/exp3_dim16_ch32_xz.png}
\includegraphics[width=0.32\textwidth]{exp3/plots/exp3_dim16_ch32_yz.png}
\includegraphics[width=0.32\textwidth]{exp3/plots/exp3_dim16_ch32_3d.png}
\caption{$k=16$, $C=32$}
\end{subfigure}
\caption{Experiment 3: intermediate dimension ($k = 16$). Left: $z$ vs.\ $\Delta(|X|)$; middle: $z$ vs.\ $\Delta(Y)$; right: both.}
\end{figure}


\section{Limitations and future work}\label{sec:limitations}

Current limitations include: batched SVD and symmetric eigensolvers lose their GPU advantage
already at block size $k = 64$ in our benchmarks (Table~\ref{tab:speedup} and
Appendix~\ref{app:heatmaps}); power iteration converges linearly in the spectral gap, fails to
separate eigenvalues $\pm\lambda$ of equal modulus, and its stopping rule controls the change of
the estimate rather than the error; automatic differentiation through the operations is not
tested; only Type~I algebras (direct sums) in finite dimensions are supported. Future work
includes support for tensor products, Cauchy-type functional calculus, distributed GPU
computing, and mixed-precision SVD. The code is open to contributions.

\section{Conclusion}\label{sec:conclusion}

We have presented \texttt{torch\_vn\_algebra}, a PyTorch library for finite-dimensional Type~I
von Neumann algebras that combines a compact tensor representation, lazy evaluation, flexible
random operator generation, and GPU-accelerated functional calculus. The library is validated
and benchmarked, and is open-source at \url{https://gitlab.com/a.hobukob/von_neumann_type_i/}
(mirror: \url{https://github.com/andre-a-no/von_neumann_type_i}).

\section*{Data availability}
All code, raw data and reproduction scripts are at the above repository.

\begin{thebibliography}{99}
\bibitem{takesaki2003} M.~Takesaki, \emph{Theory of Operator Algebras}, Springer, 2003.
\bibitem{blackadar2006} B.~Blackadar, \emph{Operator Algebras}, Springer, 2006.
\bibitem{mehta2004} M.~L.~Mehta, \emph{Random Matrices}, 3rd ed., Elsevier, 2004.
\bibitem{nielsen2010} M.~A.~Nielsen, I.~L.~Chuang, \emph{Quantum Computation and Quantum
  Information}, Cambridge University Press, 2010.
\bibitem{gardner1979} L.~T.~Gardner, ``An inequality characterizes the trace'', Canad.\ J.\
  Math.\ 31 (1979), 1322--1328.
\bibitem{novikov2015} A.~Novikov, O.~Tikhonov, ``Characterization of central elements of operator
  algebras by inequalities'', arXiv:1502.01267 (2015).
\bibitem{novikov2017} A.~Novikov, ``$L_1$-space for a positive operator affiliated with von
  Neumann algebra'', Positivity 21 (2017), 359--375.
\bibitem{petz1988} D.~Petz, J.~Zem\'anek, ``Characterizations of the trace'', Linear Algebra
  Appl.\ 111 (1988), 43--52.
\bibitem{virosztek2017} D.~Virosztek, ``Connections between centrality and local monotonicity of
  certain functions'', J.\ Math.\ Anal.\ Appl.\ 453 (2017), 221--226.
\bibitem{michelson1927} A.~A.~Michelson, \emph{Studies in Optics}, University of Chicago Press,
  1927.
\bibitem{abed2024} S.~A.~Abed, I.~A.~Nikolaeva, A.~A.~Novikov, ``Generalisation of Michelson
  contrast for operators and its properties'', Lobachevskii J.\ Math.\ 45 (2024), 3835--3848.
\bibitem{novikov2020} A.~Novikov, S.~A.~Abed, I.~Nikolaeva, ``Generalisation of Michelson contrast
  for operators'', arXiv:2010.10130 (2020).
\bibitem{mezzadri2007} F.~Mezzadri, ``How to generate random matrices from the classical compact
  groups'', Notices Amer.\ Math.\ Soc.\ 54 (2007), 592--604.
\bibitem{paszke2019} A.~Paszke et al., ``PyTorch: An imperative style, high-performance deep
  learning library'', NeurIPS 2019.
\bibitem{harris2020} C.~R.~Harris et al., ``Array programming with NumPy'', Nature 585 (2020),
  357--362.
\bibitem{virtanen2020} P.~Virtanen et al., ``SciPy 1.0: fundamental algorithms for scientific
  computing in Python'', Nat.\ Methods 17 (2020), 261--272.
\bibitem{johansson2013} J.~R.~Johansson, P.~D.~Nation, F.~Nori, ``QuTiP 2: A Python framework for
  the dynamics of open quantum systems'', Comput.\ Phys.\ Commun.\ 184 (2013), 1234--1240.
\bibitem{qiskit2019} Qiskit contributors, ``Qiskit: An open-source framework for quantum
  computing'', 2019, \url{https://qiskit.org}.
\end{thebibliography}

\appendix
\section{Complete set of speedup heatmaps}\label{app:heatmaps}

\texttt{scripts/benchmark.py} generates a heatmap for every operation.
Figures~\ref{fig:hm-first}--\ref{fig:hm-last} show the mean GPU/CPU speedup as a function of the
block size $k$ (vertical axis) and the number of channels $C$ (horizontal axis); brighter
colours (viridis colour map) mean larger speedup, values below $1$ mean that the CPU is faster.
In the stored run, \texttt{lambda\_max} was computed by exact diagonalisation of $AA^T$,
\texttt{entropy} from normalised singular values and \texttt{michelson} from the matrix entries
of Gaussian matrices; the current script evaluates them on positive definite matrices from
eigenvalues (see \texttt{results/benchmark/README.md}).

\begin{figure}[p]\centering
\includegraphics[width=\textwidth]{speedup_heatmap_trace.png}
\caption{Speedup heatmap: trace.}\label{fig:hm-first}
\end{figure}
\begin{figure}[p]\centering
\includegraphics[width=\textwidth]{speedup_heatmap_addition.png}
\caption{Speedup heatmap: addition.}
\end{figure}
\begin{figure}[p]\centering
\includegraphics[width=\textwidth]{speedup_heatmap_multiplication.png}
\caption{Speedup heatmap: multiplication.}
\end{figure}
\begin{figure}[p]\centering
\includegraphics[width=\textwidth]{speedup_heatmap_lambda_max.png}
\caption{Speedup heatmap: largest eigenvalue.}
\end{figure}
\begin{figure}[p]\centering
\includegraphics[width=\textwidth]{speedup_heatmap_svd_abs.png}
\caption{Speedup heatmap: singular values (\texttt{svd\_abs}).}
\end{figure}
\begin{figure}[p]\centering
\includegraphics[width=\textwidth]{speedup_heatmap_inverse.png}
\caption{Speedup heatmap: inverse.}
\end{figure}
\begin{figure}[p]\centering
\includegraphics[width=\textwidth]{speedup_heatmap_frobenius_norm.png}
\caption{Speedup heatmap: Frobenius norm.}
\end{figure}
\begin{figure}[p]\centering
\includegraphics[width=\textwidth]{speedup_heatmap_trace_norm.png}
\caption{Speedup heatmap: trace (nuclear) norm.}
\end{figure}
\begin{figure}[p]\centering
\includegraphics[width=\textwidth]{speedup_heatmap_entropy.png}
\caption{Speedup heatmap: entropy.}
\end{figure}
\begin{figure}[p]\centering
\includegraphics[width=\textwidth]{speedup_heatmap_michelson.png}
\caption{Speedup heatmap: Michelson contrast.}\label{fig:hm-last}
\end{figure}

\end{document}
